% Standalone problem document, generated from the adjacent README.md.
% Compile with XeLaTeX. Mathematical content is maintained in Markdown.
\documentclass[11pt,a4paper]{article}
\usepackage[fontsize=10.5pt]{scrextend}
\usepackage[margin=22mm,headheight=15pt,headsep=9mm,footskip=11mm]{geometry}
\usepackage{fontspec}
\setmainfont{texgyrepagella}[Extension=.otf,UprightFont=*-regular,BoldFont=*-bold,ItalicFont=*-italic,BoldItalicFont=*-bolditalic]
\setsansfont{texgyreheros}[Extension=.otf,UprightFont=*-regular,BoldFont=*-bold,ItalicFont=*-italic,BoldItalicFont=*-bolditalic]
\setmonofont{texgyrecursor}[Extension=.otf,UprightFont=*-regular,BoldFont=*-bold,ItalicFont=*-italic,BoldItalicFont=*-bolditalic,Scale=MatchLowercase]
\usepackage{amsmath,amssymb}
\usepackage{unicode-math}
\setmathfont{texgyrepagella-math.otf}
\usepackage{xcolor,microtype,enumitem,needspace,fancyhdr}
\usepackage{bookmark}
\definecolor{ink}{HTML}{17344B}
\definecolor{accent}{HTML}{197D80}
\definecolor{muted}{HTML}{596773}
\hypersetup{colorlinks=true,linkcolor=accent,urlcolor=accent,pdftitle={MF-17 - Optimal
uniform growth after inversion of an exponentially stable
generator},pdfauthor={Open Problems in Numerical Linear Algebra}}
\urlstyle{same}
\setlength{\parindent}{0pt}
\setlength{\parskip}{5pt plus 1pt minus 1pt}
\linespread{1.04}
\setlength{\emergencystretch}{3em}
\widowpenalty=10000
\clubpenalty=10000
\displaywidowpenalty=10000
\allowdisplaybreaks[1]
\setlist{nosep,leftmargin=1.5em}
\providecommand{\tightlist}{\setlength{\itemsep}{0pt}\setlength{\parskip}{0pt}}
\providecommand{\pandocbounded}[1]{#1}
\pagestyle{fancy}
\fancyhf{}
\fancyhead[L]{\sffamily\footnotesize\color{muted}OPEN PROBLEMS IN NUMERICAL LINEAR ALGEBRA}
\fancyhead[R]{\sffamily\bfseries\footnotesize\color{accent}MF-17}
\fancyfoot[L]{\parbox[b]{0.9\textwidth}{\sffamily\fontsize{8}{10}\selectfont\color{muted}Editorial ratings; a bounded search cannot certify that a problem remains unsolved.\\Literature check: 2026-09-08}}
\fancyfoot[R]{\sffamily\footnotesize\color{muted}\thepage}
\renewcommand{\headrulewidth}{0.3pt}
\renewcommand{\headrule}{\hbox to\headwidth{\color{accent}\leaders\hrule height \headrulewidth\hfill}}
\makeatletter
\renewcommand{\section}{\Needspace{5\baselineskip}\@startsection{section}{1}{0pt}{12pt plus 2pt minus 2pt}{4pt}{\sffamily\bfseries\large\color{ink}}}
\renewcommand{\subsection}{\Needspace{4\baselineskip}\@startsection{subsection}{2}{0pt}{9pt plus 1pt minus 1pt}{3pt}{\sffamily\bfseries\normalsize\color{ink}}}
\makeatother
\setcounter{secnumdepth}{0}
\begin{document}
{\sffamily\small\color{muted}Matrix functions and stability\par}
\vspace{3pt}
{\raggedright\hyphenpenalty=10000\exhyphenpenalty=10000\sffamily\bfseries\fontsize{21}{24}\selectfont\color{ink}Optimal
uniform growth after inversion of an exponentially stable generator\par}
\vspace{7pt}
{\sffamily\small\textbf{Difficulty:} hard\quad\textbf{Importance:} interesting
to the community\par}
\vspace{4pt}
{\color{accent}\hrule height 0.8pt}
\vspace{6pt}
\textbf{Provenance:} explicit fixed-bound formalization of the source's
growth question.

\textbf{Status:} no sharp growth order found in screening on 2026-09-08.

For each fixed real number \(M>1\), let \(\mathcal A_M\) consist of all
pairs \((H,A)\) such that \(H\) is a complex Hilbert space and \(A\) is
the generator of a strongly continuous semigroup \(T(s)\) on \(H\)
satisfying \[
\|T(s)\|_{\mathcal L(H)}\leq M e^{-s}
\qquad(s\geq0).
\] Here \(A\) may be unbounded and is understood on its dense domain.
Exponential stability implies that \(0\) belongs to the resolvent set of
\(A\), so \(A^{-1}\) is bounded on \(H\) and its exponential is defined
by the operator-norm convergent power series.

Determine, for each fixed \(M>1\), the asymptotic order as
\(t\to\infty\) of \[
G_M(t)=
\sup_{(H,A)\in\mathcal A_M}
\|\exp(tA^{-1})\|_{\mathcal L(H)}.
\] An asymptotic order means a rate \(g_M(t)\) with matching positive
upper and lower multiplicative bounds for all sufficiently large \(t\);
their constants may depend on \(M\), but must be uniform over
\((H,A)\in\mathcal A_M\).

The source's Problem 1.3 asks for the optimal inverse-semigroup growth
for exponentially stable Hilbert-space semigroups. The envelope above
makes the dependence on the original semigroup bound explicit. This
formalization is supported by uniform upper estimates and examples in
the same normalized class, rather than attributed verbatim to the
source. The problem concerns operator stability and rational
time-stepping analysis, extending finite-dimensional matrix exponential
questions.

\subsection{References}\label{references}

\begin{enumerate}
\def\labelenumi{\arabic{enumi}.}
\tightlist
\item
  E. Lorist, M. Meyries, and M. Veraar, \emph{A solution to the inverse
  generator problem and related questions}, arXiv:2608.06272v3 (2026),
  Problem 1.3, Theorem 1.2(2), and Section 1.3 on numerical stability.
  \href{https://arxiv.org/html/2608.06272v3}{Paper}.
\item
  C. J. K. Batty, A. Gomilko, and Y. Tomilov, \emph{A Besov algebra
  calculus for generators of operator semigroups and related
  norm-estimates}, Mathematische Annalen 379 (2021), pp.~23--93,
  Corollary 5.7: explicit constants depending only on the semigroup
  bound and decay rate. Their generator notation is \(-A\).
  \href{https://link.springer.com/article/10.1007/s00208-019-01924-2}{Published
  paper}.
\item
  H. Zwart, \emph{Growth Estimates for \(\exp(A^{-1}t)\) on a Hilbert
  Space}, Semigroup Forum 74 (2007), pp.~487--494, Theorem 2.2 and
  Corollary 2.3: the logarithmic upper bound.
  \href{https://link.springer.com/article/10.1007/s00233-006-0679-1}{Published
  paper}.
\end{enumerate}

\subsection{Status check --- 2026-09-08}\label{status-check-2026-09-08}

checked the current arXiv history, with latest version v3 dated
2026-08-26 and no withdrawal, and searched ``inverse generator Problem
1.3'' and ``inverse generator growth 2026 logarithmic''. Corollary 5.7
of reference 2 gives, in the present sign convention and at decay rate
one, \[
G_M(t)\leq2M^2\bigl(2-e^{-1}+e^{-1}\log t\bigr),
\qquad t>1.
\] Theorem 1.2(2) of reference 1 gives examples with a doubly
logarithmic lower bound raised to a positive exponent, while the
original semigroup bound is \(1+C\alpha/(1-\alpha)\) for an absolute
\(C\) and \(0<\alpha<1\). Thus its parameters allow examples within each
fixed \(M>1\) class. These estimates do not determine the matching
asymptotic order. The older yes/no inverse-generator problem was
resolved negatively in the 2026 paper and is not listed as open here.
\end{document}
